\documentclass[11pt,a4paper,sans]{moderncv}
\pagestyle{empty}
\moderncvstyle{classic} % styles: casual, classic, banking, oldstyle, fancy
\moderncvcolor{blue}    % colors: blue, orange, green, red, purple, grey, black

\usepackage{enumitem}
\usepackage{multicol}
\setlength{\columnsep}{3em}
\usepackage[utf8]{inputenc}
\usepackage[T1]{fontenc}
\usepackage[english]{babel}
\usepackage[a4paper, top=1.5cm, bottom=1.5cm, left=1.5cm, right=1.5cm]{geometry}

\name{Raphaël}{Jontef}
\photo[64pt][0.4pt]{../../photo.jpg}
\phone[mobile]{+33 7 68 14 51 09}
\email{raphaeljontef@hotmail.com}

\date{août 2026} % nouveau, à voir si ça rend bien

\begin{document}
\makecvtitle

\section{Objectif}
Postuler à un stage en entreprise de juin 2027 à septembre 2027.

\section{Formation}
\cventry{Paris 20\textsuperscript{e} arr.}{sept. 2020 -- juil. 2023}{Lycée Maurice Ravel}{}{}{Spécialités NSI et Mathématiques, mention "bien" au baccalauréat général et technologique.}
\vspace{1em}
\cventry{Paris 8\textsuperscript{e} arr.}{sept. 2023 -- juil. 2025}{Lycée Fénelon Sainte-Marie}{CPGE}{}{MP2I puis MPI (Mathématiques, Physique et Informatique)}
\vspace{1em}
\cventry{Bordeaux, Talence}{sept. 2025 -- 2028}{Enseirb-Matmeca}{diplôme d'ingénieur}{2ème année}{Filière Informatique}.

\section{Expériences professionnelles et projets personnels}
\vspace{0.5em}
\cvitem{Enseignement}{Tutorat présenciel en sciences pour des élèves de collège et lycée, chez \href{https://www.alveusclub.com/}{Alveus}. Environ 70 heures de cours dispensées de novembre 2025 à aujourd'hui, principalement en physique et en mathématiques.}
\cvitem{Prog. embarquée}{Conception intégrale (boitier, firmware sur ESP32) d'un ordinateur de bord pour vélo, avec indication GPS, heure, commande via Bluetooth depuis un smartphone.}
\vspace{0.5em}
\cvitem{Sensibilisation}{Animation en juin 2026 d'un stage destiné à des élèves de seconde "Apprendre à discuter avec l'IA".}
\vspace{0.5em}
\cvitem{Autre}{Travail d'Initiative Personnelle Encadrée (TIPE présenté en juillet 2025) en CPGE, sur le thème de la reconnaissance de caractère par ordinateur pour retrouver du code source \LaTeX\,à partir d'une image.}

\section{Compétences informatiques}
\begin{itemize}[left=0pt, label=\textbullet]
\begin{multicols}{2}
\item paradigme impératif avec \texttt{C} et outils associés, \texttt{Python}, ou fonctionnel (\texttt{OCaml}) avec bases en orienté-objet;
\item Utilisation de \href{https://fr.wikipedia.org/wiki/NixOS}{NixOS} sur mes machines de travail et mon serveur domestique.
    \item écriture de documents en \LaTeX;
    \item utilisation quotidienne de \texttt{git} et de l'environnement Linux en général.
\end{multicols}
\end{itemize}

\noindent
\begin{minipage}[t]{0.48\textwidth}  % gauche
\section{Compétences humaines}
\begin{itemize}[left=0pt, label=\textbullet]
    \item Bonne communication, notamment à l'écrit;
    \item Pédagogie, recherche de méthodes d'explication adpatées à l'enseignement ou à la vulgarisation;
    \item Sens affûté du détail;
    \item Importante volonté de bien faire.
\end{itemize}
\end{minipage}
\hfill
\begin{minipage}[t]{0.48\textwidth}  % droite
\section{Passe-temps}
\begin{itemize}[left=0pt, label=\textbullet]
    \item Cuisine et patisserie;
    \item cyclisme, aussi bien par nécessité que par loisir;
    \item mécanique vélo, notamment entretien et réparation;
    \item grand intérêt à comprendre le fonctionnement de tous types de systèmes mécaniques, électriques ou électroniques;
    \item passion pour les mathématiques appliquées, surtout venant de l'algèbre.
\end{itemize}
\end{minipage}

\end{document}
